\chapter{Pochodne cząstkowe i kierunkowe}
\label{ma2:ch:directional-derivatives}

W rozdziale \ref{ma2:ch:multivariable-continuity}
badaliśmy granice przy dowolnym sposobie zbliżania się
do punktu. Teraz rozpatrujemy ilorazy różnicowe wzdłuż
prostych. Korzystamy z ogólnych definicji i przykładów
analizy funkcjonalnej, rozdział 14 (dalej AF 14),
oraz z rachunku pochodnych jednej zmiennej w AM I 20.

\section{Pochodne kierunkowe i cząstkowe}

Niech $n,m$ będą dodatnimi liczbami całkowitymi,
$U\subseteq\R^n$ zbiorem otwartym, $f:U\to\R^m$
i $p\in U$. Dla $h\in\R^n$ istnieje $\delta>0$ takie,
że $p+th\in U$ dla $|t|<\delta$: jeśli
$B(p,r)\subseteq U$, dla $h\ne0$ można wziąć
$\delta=r/\|h\|_2$, a dla $h=0$ dowolne $\delta>0$.

\begin{definicja}[Pochodna kierunkowa i cząstkowa]
\label{ma2:def:directional}
Dwustronną pochodną kierunkową odwzorowania $f$ w punkcie
$p$ w kierunku $h$ nazywamy granicę
\[
 \partial_h f(p)=\lim_{\substack{t\to0\\t\in\R,\ t\ne0}}
 \frac{f(p+th)-f(p)}{t},
\]
o ile istnieje ona w $\R^m$. Jest to definicja AF 14.1.1
w przestrzeniach euklidesowych.
Dla $j=1,\ldots,n$, oznaczając przez $e_j$ $j$-ty wektor
standardowej bazy $\R^n$, definiujemy pochodną cząstkową
\[
 \partial_j f(p)=\frac{\partial f}{\partial x_j}(p)
 =\partial_{e_j}f(p).
\]
Pochodną jednostronną w kierunku $h$ otrzymujemy,
zastępując $t\to0$ przez $t\to0^+$.
\end{definicja}

Nie wymagamy $\|h\|_2=1$. Gdy $h=0$, iloraz jest zerowy
i $\partial_0 f(p)=0$. Jeśli $\gamma_h(t)=f(p+th)$,
to $\partial_h f(p)=\gamma_h'(0)$ z definicji pochodnej
krzywej. W szczególności pochodna cząstkowa jest
pochodną funkcji jednej zmiennej, w której pozostałe
współrzędne są ustalone.

\begin{stwierdzenie}[Rachunek współrzędnych]
\label{ma2:prop:directional-components}
Niech $f=(f_1,\ldots,f_m)$. Pochodna $\partial_h f(p)$
istnieje wtedy i tylko wtedy, gdy istnieje każda
$\partial_h f_i(p)$; wówczas
\[
 \partial_h f(p)=(\partial_h f_1(p),\ldots,
 \partial_h f_m(p)).
\]
Jeśli $\partial_h f(p)$ istnieje, to dla $\lambda\in\R$
istnieje $\partial_{\lambda h} f(p)$ i
$\partial_{\lambda h}f(p)=\lambda\partial_h f(p)$.
\end{stwierdzenie}
\begin{proof}
Pierwsza teza wynika z zastosowania kryterium
\ref{ma2:thm:multi-limit} do odwzorowania jednej
zmiennej $t\mapsto(f(p+th)-f(p))/t$.
Dla $\lambda=0$ obie strony drugiej równości są zerowe.
Dla $\lambda\ne0$ podstawiamy $s=\lambda t$:
\[
 \frac{f(p+t\lambda h)-f(p)}{t}
 =\lambda\frac{f(p+sh)-f(p)}s.
\]
Przy $t\to0$ parametr $s$ dąży do zera dwustronnie,
również dla $\lambda<0$. Ciągłość mnożenia przez skalar
(twierdzenie \ref{ma2:thm:metric}) daje tezę.
\end{proof}

\begin{przyklad}
Dla $f(x,y)=x^2y+\sin y$ mamy
\[
 \partial_1 f(x,y)=2xy,\qquad
 \partial_2 f(x,y)=x^2+\cos y.
\]
Pierwszy wzór uzyskujemy, traktując $y$ jako stałą,
a drugi, traktując $x$ jako stałą; stosujemy AM I,
20.3.1--2 i 20.5.2.
Dla $F(x,y)=(x^2+y,xy)$ kryterium współrzędnych
daje $\partial_1F(x,y)=(2x,y)$ i
$\partial_2F(x,y)=(1,x)$.
\end{przyklad}

\section{Interpretacja geometryczna}

Dla funkcji skalarnej $f:U\to\R$ oraz jednostkowego
$u$ pochodna $\partial_u f(p)$ opisuje zmianę wartości
na jednostkę przesunięcia wzdłuż prostej $p+tu$.
Krzywa przekroju wykresu ma postać
$t\mapsto(p+tu,f(p+tu))\in\R^{n+1}$.
Jeśli pochodna istnieje, jej wektor styczny w $t=0$
to $(u,\partial_u f(p))$, ponieważ iloraz różnicowy
pierwszej części jest stale równy $u$, a ostatniej
dąży do $\partial_u f(p)$.
Nie wyciągamy stąd wniosku o istnieniu wspólnej
płaszczyzny stycznej do całego wykresu; potrzebna
będzie różniczkowalność z następnego rozdziału.

\begin{figure}[htbp]
\centering
\begin{tikzpicture}[>=Stealth,x=1.35cm,y=1.1cm]
 \draw[->] (-1.65,0)--(2.05,0) node[right] {$t$};
 \draw[->] (0,-.25)--(0,3.45) node[above] {$f(p+tu)$};
 \draw[blue!70!black,thick,domain=-1.3:1.55,samples=65]
 plot (\x,{1.15+.65*\x+.36*\x*\x});
 \draw[red!65!black,thick] (-1.3,.305)--(1.45,2.0925)
 node[right] {styczna};
 \coordinate (P) at (0,1.15);
 \coordinate (Q) at (1.2,2.4484);
 \draw[gray,densely dashed] (P)--(1.2,1.15)--(Q);
 \draw[gray,thick] (P)--(Q);
 \fill (P) circle (1.5pt) node[left] {$f(p)$};
 \fill (Q) circle (1.5pt) node[above left] {$f(p+tu)$};
 \node[below] at (.6,1.15) {$t$};
 \node[right] at (1.2,1.78) {$\Delta f$};
 \node[blue!70!black,anchor=west] at (-1.5,2.7)
 {$\gamma_u(t)=f(p+tu)$};
\end{tikzpicture}
\caption{Pochodna wzdłuż kierunku jednostkowego jako
nachylenie stycznej do przekroju. Iloraz $\Delta f/t$
dąży do $\partial_u f(p)$.}
\label{ma2:fig:directional-section}
\end{figure}

\section{Reguły rachunkowe}

\begin{twierdzenie}
\label{ma2:thm:directional-rules}
Niech $f,g:U\to\R$ mają pochodne w kierunku $h$
w punkcie $p$. Dla $\alpha,\beta\in\R$ zachodzą
\begin{align*}
 \partial_h(\alpha f+\beta g)(p)
 &=\alpha\partial_hf(p)+\beta\partial_hg(p),\\
 \partial_h(fg)(p)
 &=\partial_hf(p)g(p)+f(p)\partial_hg(p),\\
 \partial_h(f/g)(p)
 &=\frac{\partial_hf(p)g(p)-f(p)\partial_hg(p)}
 {g(p)^2},\qquad g(p)\ne0.
\end{align*}
Jeśli $J\subseteq\R$ jest otwartym przedziałem,
$f(U)\subseteq J$, a $\varphi:J\to\R$ jest
różniczkowalna w $f(p)$, to
\[
 \partial_h(\varphi\circ f)(p)
 =\varphi'(f(p))\partial_h f(p).
\]
\end{twierdzenie}
\begin{proof}
Stosujemy twierdzenie AM I 20.3.1 do funkcji
$t\mapsto f(p+th)$ i $t\mapsto g(p+th)$ w $t=0$.
Istnienie ich pochodnych zapewnia ciągłość w zerze;
jeśli $g(p)\ne0$, mianownik pozostaje niezerowy
dla dostatecznie małych $|t|$. Ostatnia równość
wynika z jednowymiarowej reguły łańcuchowej
(AM I 20.4.1) dla $t\mapsto\varphi(f(p+th))$.
Nie jest to ogólna reguła łańcuchowa dla dwóch
odwzorowań wielowymiarowych.
\end{proof}

\begin{przyklad}
Niech $f(x,y)=\exp(xy)$. Dla $p=(a,b)$ i $h=(u,v)$
mamy $(a+tu)(b+tv)=ab+t(av+bu)+t^2uv$.
Z AM I 20.5.1 i twierdzenia \ref{ma2:thm:directional-rules}
otrzymujemy
\[
 \partial_h f(a,b)=e^{ab}(bu+av).
\]
W szczególności $\partial_1 f(a,b)=be^{ab}$
i $\partial_2 f(a,b)=ae^{ab}$.
\end{przyklad}

\section{Różniczkowalność Gâteaux i gradient}

\begin{definicja}
\label{ma2:def:gateaux}
Odwzorowanie $f:U\to\R^m$ jest różniczkowalne
w sensie Gâteaux w $p$, jeżeli istnieje operator
liniowy $A:\R^n\to\R^m$ taki, że
$\partial_hf(p)=Ah$ dla każdego $h\in\R^n$.
Oznaczamy $D_Gf(p)=A$. Jest to definicja AF 14.1.2:
w skończonym wymiarze operator liniowy jest ciągły,
gdyż, pisząc $h=\sum_jh_je_j$, mamy
\[
 \|Ah\|_2\le \sum_j|h_j|\|Ae_j\|_2
 \le \|h\|_2\sum_j\|Ae_j\|_2.
\]
\end{definicja}

\begin{stwierdzenie}
\label{ma2:prop:gateaux-matrix}
Jeśli $D_Gf(p)$ istnieje, jego macierz w standardowych
bazach ma postać
\[
 J_f(p)=(\partial_j f_i(p))_
 {1\le i\le m,\ 1\le j\le n}.
\]
W szczególności
$\partial_h f(p)=\sum_{j=1}^n h_j\partial_j f(p)$.
\end{stwierdzenie}
\begin{proof}
Z definicji $Ae_j=\partial_j f(p)$, więc $j$-tą
kolumną macierzy $A$ jest wektor $\partial_j f(p)$.
Liniowość daje wzór na $Ah$ i jednoznaczność $A$.
\end{proof}

Dla funkcji skalarnej różniczkowalnej w sensie Gâteaux
definiujemy gradient
\[
 \nabla f(p)=(\partial_1 f(p),\ldots,\partial_n f(p)).
\]
Wtedy $\partial_h f(p)=\langle\nabla f(p),h\rangle$
w standardowym iloczynie skalarnym, na mocy stwierdzenia
\ref{ma2:prop:gateaux-matrix}. W szczególności, gdy istnieje
pochodna Frécheta, jest to gradient z AF 14.1.24,
ponieważ pochodne Frécheta i Gâteaux są wtedy równe
(AF 14.1.5). Jeśli istnieją tylko pochodne cząstkowe,
ich wektor także bywa nazywany gradientem, ale sam
w sobie nie daje wzoru na pochodną w dowolnym kierunku.

\begin{wniosek}[Największa zmiana kierunkowa]
\label{ma2:cor:gradient-steepest}
Niech $f:U\to\R$ będzie różniczkowalna w sensie
Gâteaux w $p$ i niech $v=\nabla f(p)$.
Wówczas
\[
 \max_{\|u\|_2=1}\partial_u f(p)=\|v\|_2,\qquad
 \min_{\|u\|_2=1}\partial_u f(p)=-\|v\|_2.
\]
Jeśli $v\ne0$, maksimum osiągane jest dokładnie dla
$u=v/\|v\|_2$, a minimum dla $u=-v/\|v\|_2$.
Dla $v=0$ wszystkie pochodne kierunkowe są zerowe.
\end{wniosek}
\begin{proof}
Z nierówności Cauchy'ego--Schwarza
$|\langle v,u\rangle|\le\|v\|_2\|u\|_2=\|v\|_2$.
Podstawienie wskazanych wektorów jednostkowych daje
odpowiednie wartości. Dla $v\ne0$ równość w oszacowaniu
wymaga liniowej zależności $u,v$, a znak iloczynu
skalarnego rozstrzyga, który z dwóch wektorów otrzymujemy.
Dla $v=0$ wzór na pochodną daje zero w każdym kierunku.
\end{proof}

\section{Zależności i kontrprzykłady}

Istnienie pochodnych cząstkowych w $p$ nie oznacza
istnienia pochodnej w każdym kierunku. Dla
\[
 q(x,y)=
 \begin{cases}
 xy/(x^2+y^2),&(x,y)\ne(0,0),\\
 0,&(x,y)=(0,0)
 \end{cases}
\]
ograniczenia do osi są zerowe, więc obie pochodne
cząstkowe w zerze istnieją i wynoszą zero. Jednak
$q(t,t)=1/2$ dla $t\ne0$ i iloraz różnicowy
$q(t,t)/t=1/(2t)$ nie ma skończonej granicy dwustronnej.

Istnienie wszystkich pochodnych kierunkowych nie
oznacza addytywności względem kierunku. Wykorzystujemy
przykład AF 14.1.10 (etykieta źródłowa
\texttt{ex:nonlinear-directional14}): funkcja
$q_1(x,y)=x^3/(x^2+y^2)$ poza zerem i $q_1(0,0)=0$
ma pochodne
$\partial_{(a,b)}q_1(0,0)=a^3/(a^2+b^2)$
dla $(a,b)\ne0$, a w kierunku zerowym pochodną zero.
Dla kierunków $e_1,e_2,e_1+e_2$ wartości wynoszą
odpowiednio $1,0,1/2$, więc nie istnieje pochodna
Gâteaux. Sama funkcja jest ciągła w zerze, gdyż
$|q_1(x,y)|\le |x|$.

Nawet pochodna Gâteaux nie gwarantuje ciągłości
w punkcie. Z AF 14.1.8 (etykieta
\texttt{ex:directional-discontinuous14}) znamy
funkcję $g(x,y)=x^3y/(x^6+y^2)$ poza zerem,
$g(0,0)=0$: jej pochodna Gâteaux w zerze jest
operatorem zerowym, lecz $g(x,x^3)=1/2$ dla $x\ne0$.
Pełne obliczenia ilorazów różnicowych znajdują się
w przywołanym przykładzie. Zatem ciągłość ani tym
bardziej różniczkowalność Frécheta nie wynika z
różniczkowalności Gâteaux.

W następnym rozdziale zdefiniujemy pochodną Frécheta
jako liniowe przybliżenie całego przyrostu, z błędem
małym względem normy przyrostu. Wykorzystamy AF 14.1.3,
14.1.5--6 oraz 14.1.9, a także udowodnimy kryterium
różniczkowalności za pomocą ciągłych pochodnych
cząstkowych.
